% Dziennik promptów: Oliwier. Edytuje tylko Oliwier. Numeracja O1, O2, ...
% Nowy wpis dopisuj na końcu pliku (format: CLAUDE.md, makra: preamble.tex).

\subsection{Oliwier}
\label{app:prompty-oliwier}

\subsubsection*{Etap wyboru tematu}

Temat wybrano w rozmowie z Gemini (burza mózgów, 10 kandydatów, wybór tematu~1). Wpisy O1 i~O2
dokumentują ten etap. Dokładne wersje modelu i ustawienia nie są znane (czat webowy).

\promptentry{O1}{nieustalona (eksport rozmowy: 2026-10-28)}{Gemini (wersja nieustalona), czat webowy}{nieustalone}
  {Burza mózgów: 10 tematów dla projektu RAG, roadmapy i stacki}
\begin{lstlisting}
Przeanalizuj mi ten projekt robie z ziomalem temat 12 czyli poprostu RAG najbardziej
ciekawszy, zrob brainstorming wypisz mi z 10 tematow pod co tego raga mozna zrobic jaka
tematyke uzyc aby dobrze pokazac karmienie llma ragiem takie z ragiem i bez raga jak to
wyglada albo jakis chatbot na jakiejs stronie albo cos w tym stylu, odpal research exa
mcp i zrob mi brainstorming odnosnie tematow i wypisz do kazdego krociutka road mapke jak
do tego podejsc i stack oczywiscie. Prowadzacy polecal jakas biblioteke pythonowska do
zrobienia tego raga ale nie pamietam [...]
\end{lstlisting}
\promptresult{Lista 10 tematów z krótką roadmapą i stackiem. Wskazanie Chainlit, Streamlit i
Gradio jako bibliotek UI. Zestawienie wymagań zaliczeniowych.}

\promptentry{O2}{nieustalona (eksport rozmowy: 2026-10-28)}{Gemini (wersja nieustalona), czat webowy}{nieustalone}
  {Wybór tematu~1 (sylabusy), doprecyzowanie UI, darmowych modeli i scenariuszy odmowy}
\begin{lstlisting}
1. Uczelniany Asystent Sylabusow i Zasad Zaliczen [opis tematu z poprzedniej odpowiedzi]
Dobra ewidentnie bierzemy to, rozwin ten temat powiedz stack jak bedzie UI wygladac czy
to sam chatbot poprostu zrobimy czy jakis front w react co myslisz wg tego projektu i
wytycznych daj znac. Dodatkowo tak polecal on gradio ale daj znac ktory jest najlepszy
chainlit czy streamlit [...] czy z openroutera da sie wziac api darmowego modelu zeby
potestowac tego raga [...] Dodatkowo jak zlapac tego raga na tym ze nie bedzie chcial
odpowiedziec na wiadomosc w danym temacie [...]
\end{lstlisting}
\promptresult{Temat~1 wybrany. Rekomendacja gotowej biblioteki Pythona zamiast frontu w React,
wstępny faworyt Chainlit, darmowe modele z OpenRoutera, lokalne embeddingi, trzy scenariusze
testowe odmowy. Rekomendacje do weryfikacji w etapie wyboru stacku.}

\subsubsection*{Etap inicjalizacji projektu}

\promptentry{O3}{2026-10-28}{Claude Sonnet 5.5 (\texttt{claude-sonnet-5-5}), Claude Code}{medium}
  {Analiza pomysłu projektu, inicjalizacja repozytorium, struktura plików, szkielet raportu LaTeX}
\begin{lstlisting}
Przenalizuj mi idea calego projektu jaki tu rozpisalem (PDF: Tematy i roadmapy
projektow RAG + zrzuty wymagan przedmiotu), dodatkowo wez moja rozmowe z gemini
w ktorym dyskutowalem na temat projektu, zacznij od inicjalizacji repo daj maina
zrob go narazie publicznie, stworz README oraz przygotuj jakas sensowna strukture
plikow, cos w stylu folder raport gdzie bedzie miejsce do generowania latexa z
raportem promptow jakie uzyte z modelem; do claude.md dopisz aby w tym pliku
latex odrazu spisywal prompty uzyte z modelem i effortem, ta rozmowe z gemini
tez troche wez i dopisz juz na start ten proces jak klepnelismy temat
1. Uczelniany Asystent Sylabusow i Zasad Zaliczen [opis tematu z rozmowy: bez RAG
vs z RAG, roadmapa, stack Python/LangChain/PyPDFLoader/ChromaDB/Streamlit-Chainlit]
i opisz mi to i zaczniemy sobie od takiej wstepnej inicjalizacji projektu, a
stack w nastepnym etapie sie zajmiemy po analizie gradio/streamlit, no
przejrzymy.
\end{lstlisting}
\promptresult{Repozytorium z gałęzią \texttt{main}, README, \texttt{docs/} (wymagania, opis
projektu, plan), szkielet raportu w \texttt{raport/}, \texttt{CLAUDE.md} z regułą dziennika
promptów. Wybór stacku odłożony do następnego etapu.}

\promptentry{O4}{2026-10-28}{Claude Sonnet 5.5 (\texttt{claude-sonnet-5-5}), Claude Code}{medium}
  {Dodanie skilla latex-engineer do repo, ujednolicenie formatu LaTeX, usunięcie wpisu z bibliografii}
\begin{lstlisting}
Dodaj do repozytorium mojego autorskiego skilla latex-engineer aby collaborator tez
mogl korzystac i zeby plik sie tworzyl w podobnej formatce, w claude.md tez mozna.
Usun tez z biblografii [wklejony wpis wymagania-przedmiotu z pliku .bbl] potem jakies
sensowniejsze materialy bedziemy dodawac do biblografii, w tym co dopisales jest tylko
doslownie jaki temat mamy i wytyczne
\end{lstlisting}
\promptresult{Skill skopiowany do \texttt{.claude/skills/latex-engineer/}. Preambuła
przeorganizowana według szablonu skilla (grupy pakietów, nagłówek i stopka). Reguły LaTeX
dopisane do \texttt{CLAUDE.md}. Wpis z bibliografii usunięty, \texttt{\textbackslash printbibliography}
wyłączone do czasu pierwszych zweryfikowanych źródeł.}

\subsubsection*{Etap planowania}

\promptentry{O5}{2026-11-25}{Claude Opus 5.5 (\texttt{claude-opus-5-5}), Claude Code; subagenty Claude Sonnet 5.5 (\texttt{claude-sonnet-5-5})}{nieustalone}
  {Etap Planowania: porównanie stacku i UI, decyzje Q\&A, ADR, tablica zadań, prototyp RAG v0}
\begin{lstlisting}
[Prompt wklejony w całości; cudzysłowy, półpauzy i pauzy zamienione na znaki ASCII
na potrzeby listings.]

## Cel
Przeprowadź etap Planowania projektu "Uczelniany Asystent Sylabusów i Zasad Zaliczeń":
porównaj opcje stacku i UI, zbierz moje decyzje w krótkiej sesji Q&A, zapisz ADR oraz
2-osobową tablicę zadań w `docs/`, a dopiero po moim zatwierdzeniu zainicjalizuj minimalny,
działający prototyp RAG v0.

## Kontekst
- Repozytorium jest po inicjalizacji (etap Start). Zasady projektu i wymagania zaliczeniowe
  znajdują się w `README.md`, `CLAUDE.md`, `docs/wymagania.md`, `docs/projekt.md` oraz
  `docs/plan.md`.
- Zespół jest 2-osobowy: **Oliwier** (aktualnie inicjalizuje architekturę RAG i UI) oraz
  **Ignacy** (aktualnie scrapuje i zbiera sylabusy do `data/raw/`).
- W `CLAUDE.md` obowiązuje blokada: nie dodajemy zależności ani kodu w `src/` przed
  zatwierdzeniem stacku (UI, LLM, embeddingi, wektorówka).

## Ograniczenia krytyczne
- **NIE WOLNO** instalować pakietów, tworzyć plików w `src/`, ani edytować `CLAUDE.md` przed
  zakończeniem Kroku 2 i moją wyraźną akceptacją decyzji.
- **NIE WOLNO** wysyłać surowych sylabusów ani danych osobowych prowadzących do zewnętrznych
  usług; konfigurację kluczy API opieraj wyłącznie o zmienne środowiskowe (`.env` w
  `.gitignore`, szablon w `.env.example`).
- Dostarcz dokładnie to, o co proszę - bez nadmiarowych abstrakcji, zbędnych refaktorów i
  funkcji spoza zakresu v0.
- Dokumentacja i raport po polsku; kod, nazwy plików w `src/` oraz identyfikatory po
  angielsku.
- Przestrzegaj wszystkich reguł z `CLAUDE.md`: brak współautorstwa AI w commitach,
  obowiązkowy wpis w `raport/sections/prompty.tex` (sprawdź ostatni numer `N`) oraz test
  kompilacji `cd raport && latexmk -pdf -outdir=build main.tex`.

## Strategia sesji i delegowanie (Opus 5.5 + Sonnet 5.5)
Działasz jako orkiestrator (**Opus 5.5**). Na starcie uruchom równolegle dokładnie dwa
niezależne subagenty na modelu **Sonnet 5.5** (nie uruchamiaj dodatkowego
subagenta-weryfikatora):

1. **Subagent 1 (Sonnet 5.5 - Audyt repo i wymagań):**
   Przeczytaj `README.md`, `CLAUDE.md`, `docs/wymagania.md`, `docs/projekt.md`,
   `docs/plan.md` i końcówkę `raport/sections/prompty.tex`. Zwróć zwięzły raport
   zawierający dla każdego punktu konkrety:
   - listę wszystkich otwartych decyzji z `docs/plan.md`,
   - wymagania zaliczeniowe bezpośrednio wpływające na kod i architekturę (tryb bez RAG vs
     z RAG, cytowanie, odmowa przy braku danych, min. 3 wersje promptu, logowanie ustawień
     pod odtwarzalność, równy podział pracy),
   - aktualny stan katalogów `data/`, `src/`, `prompts/`, `eval/`.

2. **Subagent 2 (Sonnet 5.5 - Research stacku i UI):**
   Przeprowadź aktualny research (użyj dostępnych narzędzi wyszukiwania; przy każdym
   fakcie, którego nie da się potwierdzić w źródle, dopisz `[niezweryfikowane]` - nie
   zgaduj nazw modeli ani limitów) w 4 obszarach. Dla **każdej** porównywanej opcji podaj
   wszystkie wskazane kryteria:
   - **UI (`Gradio` vs `Streamlit` vs `Chainlit`):** aktywność projektu, łatwość
     przełącznika "Bez RAG / Z RAG", wbudowany podgląd pobranych chunków i źródeł,
     wsparcie języka polskiego, czas wdrożenia prototypu.
   - **LLM i polityka danych:** aktualnie dostępne darmowe modele przez OpenRouter wraz z
     wyraźnym sprawdzeniem polityki logowania/trenowania na darmowym tierze (pod kątem
     zakazu wysyłania danych objętych ograniczeniami z `docs/wymagania.md`) vs darmowa
     alternatywa lokalna (np. Ollama) lub provider bez logowania.
   - **Embeddingi dla języka polskiego + baza wektorowa:** 2-3 konkretne lokalne modele
     `sentence-transformers` mocne w języku polskim (wymiar, rozmiar, wymagany RAM,
     licencja) oraz weryfikacja `ChromaDB` vs lżejsza lokalna alternatywa.
   - **Ekstrakcja PDF/DOCX z tabelami:** porównanie `PyMuPDF`, `pdfplumber`, `Docling` i
     `pypdf` pod polskie sylabusy z tabelami ocen, godzin i efektów uczenia się.

## Przebieg pracy i bramka decyzyjna

### Krok 1 - Zbieranie danych
Uruchom oba subagenty Sonnet 5.5 równolegle i zsyntetyzuj ich wyniki.

### Krok 2 - Propozycja, podział pracy i pytania (OBOWIĄZKOWY STOP)
Przedstaw zwięźle:
1. **Tabelę porównawczą UI** + Twoją rekomendację całego stacku (UI, LLM, embeddingi,
   wektorówka, parser) z krótkim uzasadnieniem, założeniami i ryzykami (szczególnie
   prywatność/odtwarzalność).
2. **Propozycję podziału zadań 2-osobowego (`docs/tasks.md`)** dla Oliwiera i Ignacego z
   rozłącznymi katalogami/gałęziami git (uwzględniając, że Ignacy scrapuje sylabusy, a
   Oliwier buduje rdzeń RAG) oraz kontraktem danych dla `data/raw/` (format, nazewnictwo
   plików, usunięcie/anonimizacja PII prowadzących).
3. **Maksymalnie 5 numerowanych pytań decyzyjnych** zamykających otwarte punkty z
   `docs/plan.md` (przy każdym pytaniu podaj 2-3 opcje i zaznacz domyślną rekomendację,
   abym mógł odpowiedzieć skrótowo, np. `1A, 2B, 3A...`).

**ZATRZYMAJ SIĘ po zadaniu pytań.** Nie edytuj plików i nie instaluj niczego, dopóki nie
odpowiem i nie zatwierdzę decyzji.

### Krok 3 - Zapis decyzji i inicjalizacja RAG v0 (dopiero po mojej akceptacji)
Po otrzymaniu moich odpowiedzi wykonaj dokładnie tyle:
1. Utwórz `docs/adr-001-stack.md` (zatwierdzony stack, powody wyboru, odrzucone
   alternatywy, ryzyka) oraz `docs/tasks.md` (tablica zadań z właścicielem `[Oliwier]` /
   `[Ignacy]`, statusem, zakresem plików i kontraktem `data/raw/`).
2. Zaktualizuj `docs/plan.md` (odhacz zamknięte decyzje, zmień status etapu) oraz
   `CLAUDE.md` (usuń blokadę stacku, wpisz zatwierdzony stack i krótką instrukcję: "Przed
   pracą sprawdź `docs/tasks.md` i trzymaj się zakresu przypisanego właścicielowi").
3. Zainicjalizuj minimalny, uruchamialny szkielet prototypu RAG v0 w `src/`:
   - moduł ingestu i prostego chunkingu z metadanymi (`source_file`, `section`/`page`),
   - indeksowanie lokalnymi embeddingami do wybranej bazy wektorowej,
   - moduł generowania odpowiedzi obsługujący przełącznik **Bez RAG** vs **Z RAG**
     (wczytujący prompt z `prompts/v1.md`, zapisujący użyty model/temperaturę i zwracający
     listę pobranych chunków ze źródłem),
   - minimalne UI w wybranej bibliotece z przełącznikiem trybu i podglądem źródeł,
   - `requirements.txt` (lub `pyproject.toml`) + `.env.example` + uzupełnienie
     `.gitignore` o `.env` i lokalny katalog bazy wektorowej.
   - Jeśli `data/raw/` jest puste (Ignacy jeszcze scrapuje), utwórz wyłącznie 1 krótki,
     syntetyczny plik sylabusa w `data/sample/` (bez prawdziwych danych osobowych) do
     weryfikacji dymnej.
4. Zweryfikuj działanie (test importów / ingest na `data/sample/` lub obecnych plikach /
   start UI).
5. Dopisz wpis z tej sesji do `raport/sections/prompty.tex` (zgodnie z formatem z
   `CLAUDE.md`) i sprawdź kompilację `cd raport && latexmk -pdf -outdir=build main.tex`.

## Zakres plików
- **Dozwolone (dopiero w Kroku 3):** `docs/adr-001-stack.md`, `docs/tasks.md`,
  `docs/plan.md`, `CLAUDE.md`, `src/**`, `prompts/v1.md`, `data/sample/**`,
  `requirements.txt` / `pyproject.toml`, `.env.example`, `.gitignore`,
  `raport/sections/prompty.tex`.
- **Zakaz edycji:** `docs/wymagania.md`, `data/raw/**`, istniejące wpisy 1-4 w
  `raport/sections/prompty.tex`, `.claude/skills/latex-engineer/**`.

## Kryteria akceptacji
- [ ] Oba subagenty Sonnet 5.5 wykonały swoje zadania przed Krokiem 2, a przy
      niepotwierdzonych faktach widnieje `[niezweryfikowane]`.
- [ ] Agent zatrzymał się w Kroku 2 i uzyskał moją decyzję przed jakąkolwiek edycją
      plików lub instalacją pakietów.
- [ ] Powstały `docs/adr-001-stack.md` i `docs/tasks.md`, a `docs/plan.md` i `CLAUDE.md`
      są zaktualizowane.
- [ ] Prototyp v0 w `src/` uruchamia ingest (na `data/raw/` lub `data/sample/`), ma
      przełącznik **Bez RAG / Z RAG**, czyta `prompts/v1.md` i wyświetla źródła.
- [ ] Wpis został dopisany do `raport/sections/prompty.tex`, a `latexmk` skompilował
      `raport/build/main.pdf` bez błędów.

## Granice działania
- Kontynuuj samodzielnie wyłącznie odczyt i research w Kroku 1 oraz zatwierdzone edycje w
  Kroku 3.
- Zatrzymaj się i zapytaj przed: usuwaniem plików, operacjami `git commit` / `git push`,
  zmianą zakresu wykraczającą poza v0 lub gdy błąd instalacji/środowiska nie ustąpi po 2
  próbach.

## Raportowanie postępów
Raportuj status wyłącznie na przejściach między krokami oraz przy blokerach. Każde
stwierdzenie "gotowe" poprzyj wynikiem komendy lub listą zmienionych plików.

[Dopisek użytkownika w trakcie Kroku 1:]
Ten prompt do prompts oczywiscie zapisz

[Odpowiedź użytkownika na pytania Kroku 2:]
1. Gradio
2. zacznijmy od openroutera jak z setupujemy ollama lokalnei to wtedy przejdziemy a
narazie jako to do
3. docling
4. https://weii.prz.edu.pl/najnowsze-plany-studiow/plany-studiow-informatyka przejrzyj
kierumnek informatyka i Ai
5. A, ale jak bedzie lokalny to potestujemy na lokalnym sedzi

Sprawdz czy na weii da sie te plany studiow przejrzec :
https://krk.prz.edu.pl/karta.pl?mk=354&format=v1_html&C=#item-3

[Pytanie użytkownika w trakcie Kroku 3:]
NVIDIA: Nemotron 3.5 Super VL (free) NVIDIA: Nemotron 3.5 Embed 1B (free) takie modele
moze na openrouter jest sens? Takie ooglnie sa modele
[wklejona lista darmowych modeli z openrouter.ai: m.in. Nemotron 3.5 Super VL,
Nemotron 3.5 Embed 1B/8B, Nemotron 3 Ultra, Gemma 4 26B A4B, Gemma 4 31B]
\end{lstlisting}
\promptresult{Krok 1: dwa równoległe subagenty Sonnet 5.5 (audyt repo; research UI, LLM,
embeddingów, bazy wektorowej i parserów ze źródłami). Research wykluczył Chainlit
(repozytorium zarchiwizowane 2026-10-15) i wskazał, że darmowe modele OpenRouter trafiają do
dostawców, których warunki mogą dopuszczać trenowanie na danych. Krok 2: tabela UI,
rekomendacja stacku, podział zadań, 5 pytań. Decyzje użytkownika: Gradio, OpenRouter na start
(Ollama później, wbrew rekomendacji agenta), Docling, karty zajęć Informatyki (blok AI)
z \texttt{krk.prz.edu.pl}, ręczna ocena wierności. Sprawdzenie serwisu: 51 kart tylko w HTML,
każda z nazwiskami koordynatora i prowadzących, \texttt{robots.txt} z \texttt{Disallow: /}
(wymusza ręczne pobieranie). Krok 3: \texttt{docs/adr-001-stack.md}, \texttt{docs/tasks.md},
aktualizacja \texttt{docs/plan.md} i \texttt{CLAUDE.md}, prototyp v0 w \texttt{src/}
(ingest Docling, usuwanie danych osobowych, chunking po sekcjach, embeddingi
\texttt{mmlw-retrieval-roberta-large}, ChromaDB, tryby Bez RAG / Z RAG, UI Gradio ze źródłami,
log ustawień w JSONL), \texttt{prompts/v1.md}, syntetyczna karta w \texttt{data/sample/}.
Test dymny: ingest 6 chunków, trafne sekcje w wyszukiwaniu (podobieństwo 0,78--0,83),
na prawdziwej karcie usunięte 3 z 3 nazwisk, UI odpowiada HTTP 200. Niesprawdzone: samo
wywołanie LLM (brak klucza API w środowisku). Błąd w trakcie: \texttt{pkill -f} z wzorcem
obecnym we własnym poleceniu zabił powłokę agenta; serwer UI zatrzymał się poprawnie.
Dodatkowo hook \texttt{impeccable} utworzył poza zakresem plik \texttt{.impeccable/config.json}.}

\promptentry{O6}{2026-11-25}{Claude Opus 5.5 (\texttt{claude-opus-5-5}), Claude Code}{nieustalone}
  {Rozdzielenie dziennika promptów na pliki per osoba i ustalenie flow pracy}
\begin{lstlisting}
Dobra to wez zrob prompty Oliwier i prompty ignacy w tym samym miejscu bo z tego beda
caly czas konflikty wychodzic a trzeba dobry flow pracy ogarnac
\end{lstlisting}
\promptresult{Dziennik podzielony na \texttt{raport/sections/prompty/oliwier.tex} i
\texttt{ignacy.tex}; \texttt{prompty.tex} tylko dołącza oba pliki. Numeracja z prefiksem osoby
(O1, O2, \ldots{} oraz I1, I2, \ldots); dotychczasowe wpisy 1--5 przeniesione bez zmian treści
jako O1--O5 (\texttt{git mv}, historia pliku zachowana). Wpisy zeszły o poziom niżej
(\texttt{\textbackslash subsubsection*}), osoba to \texttt{\textbackslash subsection}. Zasady
w \texttt{CLAUDE.md} (wybór pliku po \texttt{git config user.name}), sekcja ,,Flow pracy''
w \texttt{docs/tasks.md}, odwołania w README i ADR zaktualizowane.}

\promptentry{O7}{2026-11-25}{Claude Opus 5.5 (\texttt{claude-opus-5-5}), Claude Code}{nieustalone}
  {Odśledzenie PDF, praca na gałęzi Oliwier-UI, test z OpenRouter (Gemma), podgląd UI}
\begin{lstlisting}
odsledz pdf i narazie nie rzucamy na main commitujemy juz na oliwier-ui branch i
przechodzimy do testow z .env openrouter gemma modelem pierwsze ten i tez chcialbym
zobaczyc juz to UI
\end{lstlisting}
\promptresult{\texttt{raport/main.pdf} usunięty z indeksu i dopisany do \texttt{.gitignore}.
Praca przeniesiona na istniejącą zdalną gałąź \texttt{Oliwier-UI} (lokalna utworzona od
\texttt{main}); nazwy gałęzi w \texttt{docs/tasks.md} poprawione na \texttt{Oliwier-UI}
i \texttt{Ignacy-data}. Test z \texttt{google/gemma-4-31b-it:free}: klucz poprawny, ale
6 z 6 zapytań odrzuconych kodem 429 (\texttt{upstream\_provider\_shared\_pool}, Google AI
Studio); \texttt{gemma-4-26b-a4b-it:free} tak samo. Po dwóch próbach test przerwany zgodnie
z zasadą z promptu O5. Diagnostycznie \texttt{nvidia/nemotron-3.5-super-vl:free} na karcie
syntetycznej: poprawna odpowiedź z cytatem [1] (11,3~s) i poprawna odmowa dla przedmiotu
spoza bazy (3,7~s); log ustawień kompletny. UI uruchomione lokalnie na porcie 7860 z modelem
Nemotron ustawionym zmienną środowiskową (bez zmiany \texttt{.env}).}
